\section{The Body That Rises}

There is a problem sitting inside the definition of 1336, and the Church has never resolved it because it was never asked to. If the souls of the just are ``truly blessed'' and ``have eternal life and rest'' \emph{before} they take up their bodies again, what is the resurrection for? On the face of it the constitution appears to have made the body a postscript to beatitude---an event that happens later to people who already have everything.

The tradition's answer to that question is the best available specimen of the whole problem this article is about: a small defined core, a serious and mostly convincing theological answer built on it, and, growing out of the answer, several centuries of detailed description that its own authors flagged as guesswork and later readers did not.

\subsection{What is actually defined}

Very little, and it is worth being exact about how little.

The Fourth Lateran Council in 1215 defined that all ``will rise again with their own bodies which they now bear, to receive according to their works.''\endnote{Fourth Lateran Council, \emph{Firmiter credimus}, DS 801: \latin{qui omnes cum suis propriis resurgent corporibus, quae nunc gestant}. Latin read 2026-07-26 at the registered patristica.net presentation of the Latin \emph{Enchiridion symbolorum}. The Catechism quotes the same clause at CCC 999.} Two things are secured by that clause and nothing else is. First, the resurrection is universal---the wicked rise too, which by itself disposes of any reading of the resurrection as a reward. Second, and this is the load-bearing word, the bodies are \emph{our own}, \latin{propria}, the ones ``we now bear.'' Whatever transformation is involved, the risen body is not a replacement issued at the door.

The Catechism, summarising, says the ``how'' is beyond us and says it flatly: ``This `how' exceeds our imagination and understanding; it is accessible only to faith.''\endnote{\emph{Catechism of the Catholic Church} 1000, official Vatican English web text, read 2026-07-26. CCC 999 quotes Lateran IV and 1 Cor 15; CCC 1001 places the resurrection ``at the last day,'' ``at the end of the world.''}

\subsection{What Scripture says, and where it stops}

Paul is asked the question directly---``But some man will say: How do the dead rise again? Or with what manner of body shall they come?''---and his answer is famous for being an answer and easy to misread as a description.\endnote{1 Cor 15:35, tracked verse text of the Douay--Rheims, read 2026-07-26; the whole passage 15:35--44 was read in the tracked artifact.} He begins by refusing the terms of the question: ``that which thou sowest is not quickened, except it die first. And that which thou sowest, thou sowest not the body that shall be: but bare grain, as of wheat, or of some of the rest. But God giveth it a body as he will.''\endnote{1 Cor 15:36--38, tracked verse text, read 2026-07-26.} A seed is genuinely continuous with the plant and gives no information about what the plant looks like. That is the analogy, and it is chosen to close the question, not to open it.

He then supplies four contrasts and stops:

\begin{studyframe}
``It is sown in corruption: it shall rise in incorruption. It is sown in dishonour: it shall rise in glory. It is sown in weakness: it shall rise in power. It is sown a natural body: it shall rise a spiritual body.''\endnote{1 Cor 15:42--44, tracked verse text, read 2026-07-26. The Douay's ``natural body'' renders \latin{corpus animale}, the body of the living soul (\latin{anima}); ``spiritual body'' renders \latin{corpus spiritale}. The contrast is not between matter and spirit---a body wholly de-materialised would not be a body---but between two principles animating the same flesh.}
\end{studyframe}

Four abstract nouns. Incorruption, glory, power, spirit. Nothing about appearance, age, dimension, or activity. And a few verses later, a text that cuts against the direction devotional imagination usually takes: ``flesh and blood cannot possess the kingdom of God: neither shall corruption possess incorruption.''\endnote{1 Cor 15:50, tracked verse text, read 2026-07-26.} Whatever the risen body is, it is not simply this body continued.

Elsewhere the New Testament adds a norm rather than a description: Christ ``will reform the body of our lowness, made like to the body of his glory.''\endnote{Phil 3:21, tracked verse text, read 2026-07-26. The Douay's ``reform'' renders \latin{reformabit}.} The pattern is the risen Christ, and what the Gospels say about the risen Christ is itself carefully strange---recognisable and not recognised, eating fish, passing through doors. And the First Epistle of John declines the whole question in one sentence: ``it hath not yet appeared what we shall be.''\endnote{1 Jn 3:2, tracked verse text, read 2026-07-26. The verse continues: ``We know that when he shall appear we shall be like to him: because we shall see him as he is''---likeness is asserted, content is not.}

\subsection{Why the body matters if the soul is already blessed}

The best Latin answer to the problem this section opened with is given, with characteristic economy, in the \emph{Summa}'s treatise on happiness---and it was written before the controversy of the 1330s, which makes it a useful control.

Aquinas asks whether the body is necessary for happiness, and answers that perfect happiness does not depend on it: the intellect needs the body only for the images by which it works, and the divine essence is not seen through images, so ``without the body the soul can be happy.'' He notes, and rejects as false ``both by authority and by reason,'' the view that ``the souls of saints, when separated from their bodies, do not attain to that Happiness until the Day of Judgment''---the very position a pope would take up sixty years later.\endnote{ST I-II, q.~4, a.~5, corpus, English Dominican translation, read 2026-07-26 at New Advent. The authority he cites is 2 Cor 5:6--8.}

But he then says something subtler, in two replies that are the real answer. The separated soul's desire ``is entirely at rest, as regards the thing desired\ldots\ But it is not wholly at rest, as regards the desirer, since it does not possess that good in every way that it would wish to possess it.'' And therefore: ``after the body has been resumed, Happiness increases not in intensity, but in extent.''\endnote{ST I-II, q.~4, a.~5, reply to objection 5, read 2026-07-26. The Latin technical contrast is \latin{non intensive sed extensive}. The preceding reply (ad~4) puts it in the first person of the soul: ``the soul desires to enjoy God in such a way that the enjoyment also may overflow into the body, as far as possible\ldots\ it would still wish the body to attain to its share.''}

That is as good an answer as the question admits. The vision does not get better. The person does. A human being is not a soul temporarily inconvenienced by matter; a soul without its body has everything it wants and is not yet everything it is. The resurrection does not add to God; it finishes the creature that is looking at him.

\subsection{Augustine, guessing, and saying so}

Almost every specific thing a Western Christian believes about the risen body can be traced to the last book of \emph{The City of God}, and reading those chapters is a strange experience, because Augustine keeps interrupting himself to say he does not know.

He argues that all will rise at the age of the prime of life, on the ground that Christ rose at the stature in which he died and that ``the world's wisest men have fixed the bloom of youth at about the age of thirty''---while immediately allowing that even if someone insists everyone rises in the form in which they died, ``we need not spend much labour in disputing with him.''\endnote{Augustine, \emph{De civitate Dei} XXII.15--16, trans.\ Marcus Dods, \emph{The City of God}, vol.~II (Edinburgh: T.\,\&\,T.~Clark, 1871), public domain; read 2026-07-26 in the source library's tracked transcription of the 1871 printing at lines 20263--20313 of the artifact. Augustine's argument in XXII.15 turns on Eph 4:13, ``the measure of the age of the fulness of Christ,'' which he takes to speak of age rather than stature; XXII.18 offers the alternative reading, that the ``perfect man'' is Christ with his whole body the Church.} He argues that women rise as women, against those who inferred from ``a perfect man'' that all would rise male: ``For my part, they seem to be wiser who make no doubt that both sexes shall rise\ldots\ the sex of woman is not a vice, but nature.''\endnote{\emph{De civitate Dei} XXII.17, Dods translation, read 2026-07-26 in the tracked transcription (lines 20314--20366). His exegetical argument is that when the Lord answered the Sadducees he denied marriages and not women, and that ``They shall not be given in marriage'' applies to females and ``Neither shall they marry'' to males---which he takes as an affirmation that the sexes remain (Matt 22:30; the Douay reads ``For in the resurrection they shall neither marry nor be married, but shall be as the angels of God in heaven'').} He argues that bodily blemishes are removed while the substance remains, that whatever has been taken from the body is restored to it, and that even a body ground to powder and scattered to the winds is not beyond the Creator---``no, not a hair of its head shall perish.''\endnote{\emph{De civitate Dei} XXII.19--21, Dods translation, read 2026-07-26 in the tracked transcription (lines 20418--20617).}

And then, in the same chapter, this:

\begin{studyframe}
``But what this spiritual body shall be, and how great its grace, I fear it were but rash to pronounce, seeing that we have as yet no experience of it.''\endnote{\emph{De civitate Dei} XXII.21, Dods translation, read 2026-07-26 in the tracked transcription (line 20599 region).}
\end{studyframe}

Eight chapters later, taking up whether the risen will see God with bodily eyes, he says it again and more strongly. He begins: ``I do not say what I see, but I say what I believe.'' He works through the question at length---Eliseus saw his servant at a distance without eyes, so the saints will hardly need eyes; but the eyes will still have their office; perhaps they will be given a power of seeing incorporeal things; perhaps instead God will be so known that we shall see him ``in ourselves, in one another, in Himself, in the new heavens and the new earth, in every created thing which shall then exist''---and he brackets the whole inquiry:

\begin{studyframe}
``But as we do not know what degree of perfection the spiritual body shall attain,---for here we speak of a matter of which we have no experience, and upon which the authority of Scripture does not definitely pronounce,---it is necessary that the words of the Book of Wisdom be illustrated in us: `The thoughts of mortal men are timid, and our forecastings uncertain.'\,''\endnote{\emph{De civitate Dei} XXII.29, Dods translation, read 2026-07-26 in the tracked transcription (lines 21214--21426, the complete chapter). The internal quotation is Wisdom 9:14. The chapter's own heading in the Dods edition is ``Of the beatific vision.'' Augustine's tentative conclusion is that ``it may very well be, and it is thoroughly credible, that we shall in the future world see the material forms of the new heavens and the new earth in such a way that we shall most distinctly recognise God everywhere present.''}
\end{studyframe}

``The authority of Scripture does not definitely pronounce.'' That sentence stands in the middle of the single most influential Western text on the risen body, and it applies, by Augustine's own reckoning, to most of what surrounds it.

Aquinas noticed. Quoting Augustine's speculation about glorified eyes as an objection, he replies in six words that ought to be printed on the flyleaf of every book about heaven: ``Augustine speaks as one inquiring, and conditionally.''\endnote{ST I, q.~12, a.~3, reply to objection 2, read 2026-07-26. Aquinas goes on to quote Augustine's own conditional phrasing back at the objector and to conclude that the glorified eyes will see God as our eyes now see the life of another---that is, indirectly, the divine presence being known by the intellect immediately upon the sight of bodily things.}

\subsection{The four gifts, and where they come from}

The high medieval schools systematised the risen body's properties into four: impassibility, subtlety, agility, and clarity, drawn out of Paul's four contrasts. In the printed \emph{Summa theologiae} they occupy questions 82 to 85 of the Supplement, and the Supplement is where a reader ought to slow down.

Thomas Aquinas did not write it. He stopped work on the \emph{Summa} in December 1273 and died the following March, leaving the third part unfinished; the Supplement was assembled after his death out of his much earlier commentary on the \emph{Sentences}, by a compiler whose identity the standard older reference discusses without settling---attributed by some to Peter of Auvergne and by others to Henry of Gorkum, with the Leonine editors rejecting both, and the most probable candidate being Reginald of Piperno, his secretary.\endnote{\emph{The Catholic Encyclopedia}, art.\ ``Summa Theologiae'' (vol.~14, New York 1912), read 2026-07-26 at New Advent: ``The work has been completed by the addition of a supplement, drawn from other writings of St.\ Thomas, attributed by some to Peter of Auvergne, by others to Henry of Gorkum. These attributions are rejected by the editors of the Leonine edition\ldots\ Mandonnet\ldots\ inclines to the very probable opinion that it was compiled by Father Reginald de Piperno.'' The composition dates for the parts Thomas did write are given in the same article. This matters for weighing the material: the Supplement is authentic Thomas in substance but early Thomas, transposed by another hand, and it is not the mature judgment of the \emph{Summa} proper.} The material is genuinely his and genuinely early, rearranged by someone else. That is not a reason to dismiss it. It is a reason not to quote it as though it were the considered teaching of the mature \emph{Summa}.

What is in it is a mixture, and the mixture is instructive. Some of it is sober. On clarity, the compiled text says the brightness of the risen body ``will result from the overflow of the soul's glory into the body,'' since ``whatever is received into anything is received not according to the mode of the source whence it flows, but according to the mode of the recipient''---so that the body will differ in brightness as the soul differs in merit, as Paul says of star and star.\endnote{ST Supplement, q.~85, a.~1, corpus, English Dominican translation, read 2026-07-26 at New Advent, citing 1 Cor 15:41. The article explicitly rejects the alternative explanation, then current, that the brightness comes from a heavenly fifth essence predominating in the body.} That is an argument from a principle, and the principle is the one that runs through the whole system: glory descends from the soul's vision, not from a cosmetic applied to matter.

Some of it is a long way out. The question on subtlety spends four articles on whether a glorified body can occupy the same place as another body; the question on agility asks whether the saints will ever actually move, and whether their movement is instantaneous.\endnote{ST Supplement, qq.~83--84, article headings read 2026-07-26 at New Advent.} And here is the detail worth having: the answer to the last question is \emph{no}, and the reason is a piece of rigorous physics. A body moving from A to B must at some time be neither wholly at A nor wholly at B; the alternatives are exhaustively enumerated and disposed of; ``it is altogether impossible for a body to pass from one place to another, unless it pass through every interval.'' The glorified body ``will never attain to the dignity of the spiritual nature, just as it will never cease to be a body.''\endnote{ST Supplement, q.~84, a.~3, corpus, read 2026-07-26 at New Advent. The article rejects two named positions before giving its own, and the argument is a reductio on the continuity of local motion.}

Even in the most speculative corner of the tradition, the discipline holds: a risen body is still a body, and the argument is run to the end rather than resolved by piety. What is missing is not rigour. What is missing is data---and the Supplement, unlike Augustine, does not stop to say so.

Further out still lie the aureoles: the additional crowns assigned to virgins, martyrs, and doctors over and above the essential reward, discussed across thirteen articles, including whether an aureole is due to the angels, to Christ, and to the body.\endnote{ST Supplement, q.~96, article headings read 2026-07-26 at New Advent; a.~1 distinguishes the \latin{aurea}, the essential reward that ``consists in the perfect union of the soul with God,'' from the \latin{aureola}, an accidental reward. The distinction between essential and accidental reward is standard and is not itself in question here; the enumeration and ranking of aureoles is a school elaboration with no definitional standing.} There is nothing scandalous in this. Distinguishing an essential from an accidental reward is a reasonable thing to do, and Scripture does speak of crowns. But the reader should be clear about what has happened between Lateran IV's ``their own bodies which they now bear'' and a determination of whether virgins outrank martyrs. A great deal of thinking has happened. No further revelation has.

\subsection{Sorting it}

Defined: the resurrection is universal; it is of our own bodies; it occurs at the last day. Taught with the whole weight of Scripture and the constant tradition: the risen body is incorruptible, glorious, powerful, and ``spiritual''; it is conformed to Christ's; personal identity survives---the Catechism says the elect ``retain, or rather find, their true identity, their own name.''\endnote{CCC 1025, official Vatican English web text, read 2026-07-26, citing Phil 1:23 and Apoc 2:17. The promise of the new name given to the conqueror, ``which no man knoweth, but he that receiveth it,'' is the Apocalypse's own way of saying that identity in glory is more particular, not less.} Common theological teaching, well-argued and not defined: that beatitude increases in extent and not in intensity when the body returns; that clarity overflows from the soul; that impassibility, subtlety, agility, and clarity are the properties of the glorified body.

Speculation, marked as such by its own best author: the age of the risen, their stature, their appearance, the fate of hair and fingernails, whether the eyes see God, how movement works, and every question of the form ``what will it be like?'' Augustine wrote most of it and disclaimed it twice. Aquinas read the disclaimer and repeated it.

And two things the tradition, so far as this article has read it, simply does not teach: that the risen will be reunited in the manner devotional writing describes---the sitting down with Abraham, Isaac, and Jacob is a real promise, but it is a promise of a feast in the kingdom, not a scene at a gate\endnote{Matt 8:11, tracked verse text, read 2026-07-26: ``many shall come from the east and the west, and shall sit down with Abraham, and Isaac and Jacob in the kingdom of heaven.'' The parable of Dives and Lazarus (Luke 16:19--31) shows the dead recognising and addressing one another, but it is a parable and its topography is the parable's; nothing in the passage is offered by the Church as a description of the state of the blessed.}---and that the risen body will be young, beautiful, and thirty. That last is Augustine's opinion, offered with a shrug, and it has had a career its author would have found startling.
